% Generated draft: requires mathematical and source review.
\documentclass[11pt]{article}
\usepackage[margin=25mm]{geometry}
\usepackage{fontspec}
\setmainfont{Latin Modern Roman}
\usepackage{amsmath,amssymb,mathtools,unicode-math}
\setmathfont{Latin Modern Math}
\usepackage{enumitem,needspace,xurl,hyperref}
\setlength{\parindent}{0pt}
\setlength{\parskip}{5pt}
\setlength{\emergencystretch}{3em}
\newcommand{\N}{\mathbb{N}}
\newcommand{\Z}{\mathbb{Z}}
\newcommand{\Q}{\mathbb{Q}}
\newcommand{\R}{\mathbb{R}}
\newcommand{\C}{\mathbb{C}}
\newcommand{\F}{\mathbb{F}}
\newcommand{\A}{\mathbb{A}}
\newcommand{\PP}{\mathbb{P}}
\newcommand{\E}{\mathbb{E}}
\newcommand{\eps}{\varepsilon}
\newcommand{\dd}{\,\mathrm{d}}
\begin{document}
\section*{UnsolvedMath ID: 30001536}
\textbf{Source problem code:} catalogue entry
\subsection{Definitions and mathematical statement}
Let \(G\) be a finite nontrivial group and let \(g\in\Z_{>0}\). Write \(G^{\mathrm{ab}}=G/[G,G]\), where \([G,G]\) is the commutator subgroup, and define
\[
d(G^{\mathrm{ab}})=\min\{\lvert S\rvert:S\subseteq G^{\mathrm{ab}},\ \langle S\rangle=G^{\mathrm{ab}}\}.
\]
In particular, the trivial group has generator number zero. A \(G\)-extension of \(\Q\) means a finite Galois extension \(K/\Q\) with \(\operatorname{Gal}(K/\Q)\cong G\). No restriction to tame ramification is stated in the supplied record.

For such an extension, let
\[
\operatorname{Ram}_{\mathrm{fin}}(K/\Q)=\{p\text{ a rational prime}:p\text{ ramifies in }K/\Q\}.
\]
The real place \(\infty\) of \(\Q\) is ramified when a place of \(K\) above it is complex; since \(K/\Q\) is Galois, this is equivalent to \(K\) being totally imaginary. Define
\[
\epsilon_\infty(K)=\begin{cases}1,&\infty\text{ ramifies in }K/\Q,\\0,&\infty\text{ does not ramify in }K/\Q,\end{cases}
\qquad
r(K/\Q)=\lvert\operatorname{Ram}_{\mathrm{fin}}(K/\Q)\rvert+\epsilon_\infty(K).
\]
Thus ramified places of the base field are counted once each, including the infinite place when it ramifies.

Under the interpretation that ``including the infinite prime'' specifies this counting convention, set
\[
\mathcal E_{G,g}=\{K/\Q:K/\Q\text{ is a }G\text{-extension and }r(K/\Q)=g\}/\cong_{\Q}.
\]
The supplied problem asks, for every finite nontrivial \(G\) and every positive integer \(g\), whether \(\mathcal E_{G,g}\) is nonempty and, when it is nonempty, how many such extensions there are. Counting up to \(\Q\)-isomorphism is an explicit convention for this draft; the source record does not specify a counting equivalence. Its proposed existence criterion is
\[
\mathcal E_{G,g}\ne\varnothing
\quad\overset{?}{\Longleftrightarrow}\quad
g\ge d(G^{\mathrm{ab}}).
\]
The equality \(r(K/\Q)=g\) means exactly \(g\) ramified places, rather than at most \(g\).

There is an unresolved alternative reading: the words ``including the infinite prime'' may require \(\infty\) itself to ramify. Under that reading, replace \(\mathcal E_{G,g}\) throughout by
\[
\mathcal E^{\infty}_{G,g}=\{[K]\in\mathcal E_{G,g}:\epsilon_\infty(K)=1\},
\]
so that exactly \(g-1\) finite rational primes ramify. These readings impose different conditions, and the supplied record does not settle which is intended. The displayed criterion therefore records the supplied conjectural wording with this qualification; it is not asserted here as a mathematically verified conjecture. In particular, the draft retains the generator number of \(G^{\mathrm{ab}}\) exactly as supplied. The associated background describes an asserted minimal-ramification equality, but does not explain its relationship to existence for every exact count \(g\) above the stated threshold.
\subsection{Short English statement}
For a finite nontrivial group \(G\) and a positive integer \(g\), ask whether there is a Galois extension of \(\Q\) with group \(G\) and exactly \(g\) ramified base places, and how many such extensions exist. The supplied criterion is \(g\ge d(G^{\mathrm{ab}})\). Whether the infinite place must ramify, or merely counts when it does, remains unresolved. The supplied open status and conjectural formulation require human mathematical and source review.
\subsection{Sources}
\begin{itemize}
\item Supplied local UnsolvedMath registry record, permanent ID \(30001536\), source problem code \texttt{catalogue entry}, titled ``Minimally Ramified Extensions with Prescribed Galois Group,'' in the Oberwolfach Reports Open Problems collection. The record gives proposed year \(2010\) and names no individual proposer.
\item Supplied report citation: \emph{Low-Dimensional Topology and Number Theory} (2011), Oberwolfach Reports, pp.\ 2101--2163, DOI \texttt{10.4171/owr/2010/35}. \url{https://doi.org/10.4171/owr/2010/35}. This is also the supplied literature-source link labeled ``original OWR report.''
\item Supplied literature-source link labeled ``Boston--Markin primary paper'': \url{https://www.labmath.uqam.ca/~annales/volumes/33-2/PDF/145-154.pdf}. The supplied background attributes the minimal-ramification conjecture to Boston--Markin; it supplies no paper title or more precise location of the formulation in this reference.
\item Supplied literature assessment dated \texttt{2026-08-21}: the record labels the problem open and states that the Boston--Markin minimal-ramification conjecture is proved in abelian and several special families, while the asserted equality for all finite groups remains open, notably for broad classes of \(p\)-groups. This is supplied dated curation, not independently verified here. The background explicitly requests independent verification of time-sensitive status and cited claims.
\item The supplied statement status is \texttt{corrected\_verified}; its accompanying verification note says only that TeX notation was enclosed in mathematical delimiters without changing mathematical content. This label does not establish mathematical correctness or independent source verification.
\end{itemize}
\end{document}
